\section{Prospective scenario validation}
\label{sec:orf-prospective-supplement}
This local protocol was frozen before optimization. It fixes Budget-L2 at
0.01 MW and 100 m\textsuperscript{3}, with exactly the previous kernel,
coefficient window, integer allowance and activation export reserve. No budget
or solver parameter is tuned on these outcomes. Eight scenarios per horizon
(24,72,168 hours) use contiguous windows of the six known clean inputs, in
round-robin order, and a fixed PRNG. Inflow receives independent hourly factors
in $[0.9,1.1]$; the two demand series receive the same hourly factor, preserving
local demand below total demand. The original generator and plant parameters
are unchanged. The 24 new exports retain $4H$ binary variables, $17H$ columns
and $28H$ rows. This tests prospective parameter transfer conditional on known
parent structure, not external validation or an observational holdout.

The 360 assignments use seeds 101,137,173 and five policies in a frozen shuffled
order, sequential fresh processes on a shared host. Common settings are Gurobi
13.0.2, four threads, 120 s, 1\% requested gap, automatic presolve,
FeasibilityTol=$10^{-6}$, IntFeasTol=$10^{-9}$, ScaleFlag=$-1$ and default
IntegralityFocus/NumericFocus. Base-tight changes only FeasibilityTol to
$10^{-9}$; L2-NF3 changes only NumericFocus to 3. No starts are supplied.
Every candidate and status is retained. All policies return points and native
completion in 72/72 assignments; native completion is not a certificate.

\begin{table}[htbp]
\centering\small
\caption{Prospective scenarios. Budgets include raw and integer-rounded points;
B/G are detected bound/1\%-completion contradictions against exact feasible
witnesses of the actual native export. No missing witness contrasts occur.}
\label{tab:orf-prospective-accuracy}
\begin{tabular}{lrrr}\toprule
Policy & Budgets & B/G & Export hypothesis\\\midrule
Base & 72/72 & 0/0 & 70/72\\
Base-tight & 72/72 & 0/0 & 72/72\\
L2 & 72/72 & 7/6 & 72/72\\
L2-NF3 & 72/72 & 0/0 & 72/72\\
Budget-L2 & 72/72 & 1/0 & 72/72\\
\bottomrule\end{tabular}\end{table}
All original-unit checks use exact rational expressions in the supplied
binary64 data and points. The nominal exported tolerance is distinguished from
the observed residual: two Base points exceed the $10^{-6}$ physical-row
hypothesis slightly but still satisfy their original budgets. No policy gains
additional primal recoveries here. The contract changes 8,448 factors over the
24 models, always activation rows (14.286\% of each model's rows). The largest
Budget/L2 factor ratio per model ranges from 8.242 to 10.694. There are no
incompatible requirements. These measurements locate the changed representation,
without identifying a proprietary solver operation as its cause.

\begin{table}[htbp]
\centering\small
\caption{Paired geometric means for Budget-L2/control. Each horizon has eight
scenarios and three seeds; all rows include adverse outcomes. Time covers the
complete fresh process, including generation, selection, export, import,
verification and persistence. Certificate acquisition is separate.}
\label{tab:orf-prospective-performance}
\begin{tabular}{lrrrr}\toprule
Horizon & $W/ W_{\rm Base}$ & $W/W_{\rm L2}$ & $T/T_{\rm Base}$ & $T/T_{\rm L2}$\\\midrule
24 & 0.517 & 0.944 & 0.945 & 1.038\\
72 & 0.416 & 1.053 & 0.874 & 1.017\\
168 & 0.349 & 0.811 & 0.613 & 0.936\\
All & 0.422 & 0.931 & 0.797 & 0.996\\
\bottomrule\end{tabular}\end{table}
Across all 72 pairs, Budget-L2/Base-tight work and time ratios are 0.170 and
0.643; Budget-L2/NF3 ratios are 0.476 and 1.031. There are only 24 scenarios,
not 72 independent models. Removing pairs with a contradiction in either arm
leaves 71 Base comparisons (work/time 0.419/0.796), 71 Base-tight
(0.168/0.641), 64 L2 (0.971/1.023), and 71 NF3 (0.472/1.028).
This descriptive screen does not certify the retained gaps. Candidate cost
deterioration reaches 1.020\% against Base/Base-tight and 0.778\% against L2.
Thus the six-known-model claim of deterioration below 0.1\% does not extend to
this prospective sample. Lower work against Base is replicated, but a broad
incremental speed advantage over L2 is not.

For new-h072-08, seed 101, Budget-L2 reports a lower bound of
$-531914.2655652602$. An exactly feasible point of that actual export has cost
$-534369.9691107974$, contradicting the bound by 2455.7035455372 units.
It does not trigger the separate 1\%-completion flag. This is an adverse outcome
under the frozen policy, not a reason to retune it.

The reference campaign attempts one unfixed global rational LP per original
model and one integer-pattern repair per representation (Base,L2,Budget-L2),
each at 30 s. For each representation the lowest-cost saved candidate selects
the pattern, with a lexical tie break; no alternative patterns are tried.
All 96 LP calls are recorded. Independent reconstruction checks 72 feasible
integer witnesses, one per native representation/model, and 17 global LP dual
lower bounds. Twelve Base references have certified relative interval width
at most 1\%; none has $L=U$. Seven SoPlex global-LP attempts terminate unknown
without usable certificates, rather than necessarily exhausting the time limit.
The complete reference campaign costs 50.90 s additional to the MIP comparisons.
These repaired references do not certify every returned floating-point point.

The archive also reports a descriptive checked-completion PAR10: use complete
time for a completed, budget-passing point with a known witness and no detected
contradiction; otherwise assign 1200 s. This additional score was not an original
prespecified endpoint. The means for Base, Base-tight, L2, NF3 and Budget-L2
are 1.860,2.431,117.623,1.095 and 17.696 s. The raw observations and exclusions
remain available. No one score replaces the separate accuracy, quality and
effort comparisons.

\paragraph{Post hoc scenario-cluster uncertainty}
The original optimization protocol did not prespecify uncertainty intervals.
The subsequent analysis plan, supplied as \texttt{CLUSTER\_ANALYSIS\_PLAN.md},
fixes all four controls, both work and complete-process time, and includes every
primary pair without filtering. For each scenario and control, take the mean
of its three paired log ratios. Resample 24 complete scenarios with replacement
100,000 times; use the same sampled scenario indices for all controls and both
metrics. The 2.5th and 97.5th percentiles of the resampled mean log ratios,
with linear interpolation followed by exponentiation, give marginal 95\%
intervals. The PRNG is NumPy PCG64, seed 2026101902, with scenarios sorted by
name. No individual seed or logging repetition is an independent cluster.
These are descriptive, conditional intervals, not simultaneous intervals,
tests of superiority, or optimality certificates. Exchangeability is an
assumption of this empirical scenario mixture; sharing six known parents is
not an independently validated source of generalization. With only 24 groups,
nominal coverage is approximate and uncertainty about other parents is excluded.

The main text gives all four aggregate comparisons. The same PRNG stream
continues through horizons 24, 72 and 168, resampling eight scenarios within
each horizon. Table~\ref{tab:orf-prospective-clusters} displays the L2 comparison;
the archive gives all controls and 96 scenario/control aggregates. A scenario
win means its three-seed mean log ratio is below $-10^{-12}$; absolute values
at most $10^{-12}$ are numerical ties. No such ties occur here.

\begin{table}[htbp]
\centering\small\setlength{\tabcolsep}{4pt}
\caption{Budget-L2/L2 by horizon: geometric means and marginal 95\% percentile
intervals from resampling eight complete scenarios. Wins count scenarios out
of eight after aggregating three seeds. All work intervals include one.}
\label{tab:orf-prospective-clusters}
\begin{tabular}{lcccc}\toprule
Hours & Work ratio [95\% interval] & Time ratio [95\% interval] & Wins $W$ & Wins $T$\\\midrule
24 & 0.944 [0.818, 1.059] & 1.038 [0.980, 1.097] & 4 & 2\\
72 & 1.053 [0.839, 1.369] & 1.017 [0.921, 1.115] & 4 & 3\\
168 & 0.811 [0.382, 1.241] & 0.936 [0.562, 1.262] & 2 & 1\\
\bottomrule\end{tabular}\end{table}
Budget-L2 uses less work in 10/24 scenarios and less complete time in 6/24
against L2. The corresponding counts against Base are 21/24 and 16/24;
against Base-tight, 24/24 and 21/24; against NF3, 24/24 and 8/24.
The 168-hour work mean of 0.811 includes only two scenario wins out of eight:
unequal improvement magnitudes explain why the geometric mean alone does not
establish consistent benefit. No horizon-specific rule is inferred.
Independent replay reconstructs all 32 intervals from native API work and
outer process timings using cluster multiplicities rather than the producer's
indexed averaging. It checks every scenario aggregation and resampling hash,
without calls to an optimizer.

\paragraph{Native logging confirmation}
The packaging audit detected that the first campaign inherited OutputFlag=0,
which suppressed native traces; its API outcomes, saved points and header logs
remain intact. A registered amendment repeats all 360 assignments with
OutputFlag=1, changing no optimization parameter, model or policy. Both campaigns
are supplied: 720 MIP calls on the same 24 scenarios and three seeds, not an
enlarged independent sample. All confirmation traces contain complete optimizer
output and are checked against the saved endpoints.
The confirmation reproduces 360/360 saved points,
360/360 work values and 360/360 endpoint sets exactly.
Every repeated point is checked rationally; the unchanged exports permit reuse
of the same witnesses without new LP acquisition. Work comparisons and adverse
bound outcomes remain available for both campaigns. Complete-process time is
reported separately for this technical replication on the shared host.
Budget-L2 time ratios are 0.787 against Base, 0.637 against Base-tight, 0.986 against L2, 1.008 against L2-NF3.
